% conclusion.tex -- Sequencer Paper (BNR-2026-009)

\section{Conclusion}

This paper investigated the design and performance of a single-operator sequencer for
the Basis Network enterprise zkEVM L2. Through comparative analysis of eight production
L2 systems, implementation of a Go prototype (${\sim}$500~LOC), and comprehensive
benchmarks across scaling, concurrency, and adversarial scenarios, we establish the
following conclusions:

\begin{enumerate}
  \item Centralized sequencing with L1-based forced inclusion is the universal
    production pattern. All eight surveyed systems use centralized sequencers; forced
    inclusion via L1 provides the censorship resistance guarantee without sacrificing
    throughput.

  \item Block production overhead is negligible relative to EVM execution. At
    5{,}000~transactions per block, the sequencer consumes $<$1~ms of the 1-second
    block interval, leaving $>$99.9\% of the time budget for EVM execution.

  \item FIFO ordering is deterministic and requires no enforcement mechanism beyond
    atomic sequence numbering at insertion time. 100\% ordering accuracy is maintained
    across all tested conditions, including concurrent access with 4 producer
    goroutines.

  \item Forced inclusion achieves single-tick latency (${\sim}$1~second) under both
    cooperative and adversarial conditions, three orders of magnitude faster than the
    24-hour deadline. 100\% of forced transactions are included within the first block
    after deadline expiry.

  \item The enterprise zero-fee model eliminates MEV as a design concern. FIFO
    ordering provides fairness without the complexity of MEV-mitigation mechanisms
    required in public L2 systems.

  \item Five architectural invariants (INV-SEQ-1 through INV-SEQ-5) are established
    as constraints for downstream pipeline stages: the Logicist (TLA+ specification),
    the Architect (Go implementation), and the Prover (Coq verification).
\end{enumerate}

\subsection{Recommendations}

For the Basis Network enterprise zkEVM L2 sequencer implementation:

\begin{enumerate}
  \item Deploy a single-operator sequencer with 1-second block intervals. The
    enterprise operator runs the sequencer as part of their L2 node, with the
    forced inclusion contract on L1 providing the censorship resistance fallback.

  \item Implement FIFO mempool with atomic sequence numbering. The mutex-based
    approach is sufficient for enterprise workloads (${\sim}$33K~tx/s under
    contention, 66$\times$ above target).

  \item Deploy an Arbitrum-style forced inclusion contract on the Basis Network L1
    with a 24-hour deadline. The L1 contract maintains a FIFO queue; the sequencer
    monitors events and includes pending forced transactions at block production time.

  \item Use hybrid batch strategy (size OR time trigger) as validated in the
    validium node research (RU-V4). Block production and batch aggregation are
    separate concerns: blocks are produced every 1~second, batches are submitted
    to L1 when size or time thresholds are met.

  \item Do not implement MEV-mitigation mechanisms (threshold encryption,
    commit-reveal, fair ordering protocols). The enterprise zero-fee model makes
    these unnecessary and their complexity would degrade sequencer performance.
\end{enumerate}

\subsection{Future Work}

Immediate next steps in the Basis Network R\&D pipeline:

\begin{itemize}
  \item \textbf{RU-L3 (State Management)}: Integrate the Poseidon-based sparse Merkle
    tree with the sequencer's block production loop, measuring the combined overhead
    of state updates during block execution.
  \item \textbf{Integration benchmarks}: Combine the EVM executor (RU-L1) with the
    sequencer to measure end-to-end block production latency including transaction
    execution.
  \item \textbf{Forced inclusion contract}: Implement and deploy the Solidity L1
    contract for the forced inclusion queue, including gas cost measurements on
    the Basis Network L1.
  \item \textbf{Network layer}: Add gRPC/JSON-RPC transaction submission and measure
    end-to-end latency including serialization and network round-trip.
  \item \textbf{WAL integration}: Add Write-Ahead Log for block persistence and
    measure the overhead on block production latency.
\end{itemize}
